%%%%%%%%%%%%%%%%%%%%%%%%%%%%%%%%%%%%%%%%%%%%%%%%%%%%%%%%%%%%
% 4.7 LOW-DIMENSIONAL TOPOLOGY
% The Composition of Advanced Mathematics
%%%%%%%%%%%%%%%%%%%%%%%%%%%%%%%%%%%%%%%%%%%%%%%%%%%%%%%%%%%%

\section{Low-Dimensional Topology}
\label{sec:low-dimensional-topology}

\prerequisite{
Point-Set Topology, Algebraic Topology, Differential Topology,
Geometric Topology, Manifold Theory, and Knot Theory.
}

\learningobjectives{
By the end of this section, the reader should be able to:
\begin{itemize}
    \item explain why topology behaves differently in low dimensions;
    \item classify compact connected surfaces;
    \item distinguish orientable and nonorientable surfaces;
    \item compute the Euler characteristic of standard surfaces;
    \item describe surfaces using polygonal identifications;
    \item explain connected sums of surfaces;
    \item state the classification theorem for compact connected surfaces;
    \item describe fundamental groups of surfaces;
    \item explain the role of curves on surfaces;
    \item define mapping class groups conceptually;
    \item describe Dehn twists;
    \item explain Heegaard splittings of \(3\)-manifolds;
    \item describe surgery presentations of \(3\)-manifolds;
    \item explain prime decomposition of \(3\)-manifolds;
    \item describe incompressible surfaces;
    \item explain the JSJ decomposition conceptually;
    \item state the eight Thurston geometries;
    \item explain the geometrization viewpoint;
    \item state the Poincar\'e conjecture;
    \item explain why dimension four is exceptional;
    \item distinguish smooth and topological \(4\)-manifold phenomena;
    \item define intersection forms conceptually;
    \item explain handle decompositions and Kirby diagrams;
    \item describe exotic smooth structures;
    \item explain how knot theory interacts with \(3\)- and \(4\)-manifold topology.
\end{itemize}
}

%%%%%%%%%%%%%%%%%%%%%%%%%%%%%%%%%%%%%%%%%%%%%%%%%%%%%%%%%%%%
\subsection{What Is Low-Dimensional Topology?}
\label{subsec:what-is-low-dimensional-topology}
%%%%%%%%%%%%%%%%%%%%%%%%%%%%%%%%%%%%%%%%%%%%%%%%%%%%%%%%%%%%

Low-dimensional topology studies manifolds and embeddings primarily in
dimensions

\[
\boxed{
2,\qquad3,\qquad4.
}
\]

These dimensions display phenomena that are substantially different from
those occurring in higher-dimensional topology.

The central objects are

\[
\boxed{
\text{surfaces},
\qquad
3\text{-manifolds},
\qquad
4\text{-manifolds}.
}
\]

Knot theory also belongs naturally to low-dimensional topology because a
classical knot is an embedding

\[
S^1\hookrightarrow S^3.
\]

\begin{conceptbox}
Low-dimensional topology investigates the topology, geometry, and
classification of manifolds in dimensions \(2\), \(3\), and \(4\), where
special geometric and algebraic structures make these dimensions unusually
rich.
\end{conceptbox}

%%%%%%%%%%%%%%%%%%%%%%%%%%%%%%%%%%%%%%%%%%%%%%%%%%%%%%%%%%%%
\subsection{Why Low Dimensions Are Special}
\label{subsec:why-low-dimensions-special}
%%%%%%%%%%%%%%%%%%%%%%%%%%%%%%%%%%%%%%%%%%%%%%%%%%%%%%%%%%%%

Dimension strongly influences topology.

In dimension \(2\), compact connected manifolds can be completely
classified.

In dimension \(3\), topology interacts deeply with geometry, knots,
fundamental groups, and hyperbolic structures.

In dimension \(4\), the relationship between topology and smooth structure
becomes exceptionally subtle.

A useful schematic picture is

\[
\boxed{
\begin{array}{ccl}
2\text{D} &:& \text{surface classification},\\[2mm]
3\text{D} &:& \text{geometrization and knot complements},\\[2mm]
4\text{D} &:& \text{exotic smooth structures and intersection forms}.
\end{array}
}
\]

%%%%%%%%%%%%%%%%%%%%%%%%%%%%%%%%%%%%%%%%%%%%%%%%%%%%%%%%%%%%
\subsection{Dimension Two: Surfaces}
\label{subsec:dimension-two-surfaces}
%%%%%%%%%%%%%%%%%%%%%%%%%%%%%%%%%%%%%%%%%%%%%%%%%%%%%%%%%%%%

A connected \(2\)-manifold is called a surface.

Locally, every point has a neighborhood homeomorphic to

\[
\R^2
\]

or, for a surface with boundary, to a half-plane near a boundary point.

Familiar examples include

\[
S^2,
\qquad
T^2,
\qquad
\mathbb{RP}^2,
\qquad
\text{Klein bottle}.
\]

%%%%%%%%%%%%%%%%%%%%%%%%%%%%%%%%%%%%%%%%%%%%%%%%%%%%%%%%%%%%
\subsection{The Sphere}
\label{subsec:low-dimensional-sphere}
%%%%%%%%%%%%%%%%%%%%%%%%%%%%%%%%%%%%%%%%%%%%%%%%%%%%%%%%%%%%

The \(2\)-sphere is

\[
\boxed{
S^2
=
\left\{
(x,y,z)\in\R^3:
x^2+y^2+z^2=1
\right\}.
}
\]

It is compact, connected, orientable, and has no boundary.

Its Euler characteristic is

\[
\boxed{
\chi(S^2)=2.
}
\]

%%%%%%%%%%%%%%%%%%%%%%%%%%%%%%%%%%%%%%%%%%%%%%%%%%%%%%%%%%%%
\subsection{The Torus}
\label{subsec:low-dimensional-torus}
%%%%%%%%%%%%%%%%%%%%%%%%%%%%%%%%%%%%%%%%%%%%%%%%%%%%%%%%%%%%

The torus is denoted

\[
\boxed{
T^2=S^1\times S^1.
}
\]

It is compact, connected, orientable, and has Euler characteristic

\[
\boxed{
\chi(T^2)=0.
}
\]

Topologically, it can be constructed by identifying opposite sides of a
square.

%%%%%%%%%%%%%%%%%%%%%%%%%%%%%%%%%%%%%%%%%%%%%%%%%%%%%%%%%%%%
\subsection{Polygonal Models of Surfaces}
\label{subsec:polygonal-models-surfaces}
%%%%%%%%%%%%%%%%%%%%%%%%%%%%%%%%%%%%%%%%%%%%%%%%%%%%%%%%%%%%

Many surfaces can be represented by polygons whose edges are identified.

For the torus, one may use the boundary word

\[
\boxed{
aba^{-1}b^{-1}.
}
\]

The square's opposite sides are identified according to the labels and
orientations.

This notation compresses the topology of the entire surface into an edge
identification pattern.

%%%%%%%%%%%%%%%%%%%%%%%%%%%%%%%%%%%%%%%%%%%%%%%%%%%%%%%%%%%%
\subsection{Projective Plane}
\label{subsec:projective-plane-low-dimensional}
%%%%%%%%%%%%%%%%%%%%%%%%%%%%%%%%%%%%%%%%%%%%%%%%%%%%%%%%%%%%

The real projective plane is denoted

\[
\boxed{
\mathbb{RP}^2.
}
\]

One polygonal representation has boundary word

\[
\boxed{
aa.
}
\]

Equivalently,

\[
\mathbb{RP}^2
\]

can be obtained from a disk by identifying antipodal points on its boundary.

It is nonorientable.

Its Euler characteristic is

\[
\boxed{
\chi(\mathbb{RP}^2)=1.
}
\]

%%%%%%%%%%%%%%%%%%%%%%%%%%%%%%%%%%%%%%%%%%%%%%%%%%%%%%%%%%%%
\subsection{The Klein Bottle}
\label{subsec:klein-bottle}
%%%%%%%%%%%%%%%%%%%%%%%%%%%%%%%%%%%%%%%%%%%%%%%%%%%%%%%%%%%%

The Klein bottle is a compact connected nonorientable surface.

A standard polygonal word is

\[
\boxed{
aba^{-1}b.
}
\]

Its Euler characteristic is

\[
\boxed{
\chi(K)=0.
}
\]

Unlike the torus, the Klein bottle is nonorientable.

%%%%%%%%%%%%%%%%%%%%%%%%%%%%%%%%%%%%%%%%%%%%%%%%%%%%%%%%%%%%
\subsection{Orientability}
\label{subsec:surface-orientability}
%%%%%%%%%%%%%%%%%%%%%%%%%%%%%%%%%%%%%%%%%%%%%%%%%%%%%%%%%%%%

Informally, a surface is orientable when a consistent choice of local
orientation can be made throughout the entire surface.

Examples of orientable surfaces include

\[
\boxed{
S^2,
\qquad
T^2.
}
\]

Examples of nonorientable surfaces include

\[
\boxed{
\mathbb{RP}^2,
\qquad
\text{Klein bottle}.
}
\]

%%%%%%%%%%%%%%%%%%%%%%%%%%%%%%%%%%%%%%%%%%%%%%%%%%%%%%%%%%%%
\subsection{Connected Sum of Surfaces}
\label{subsec:connected-sum-surfaces}
%%%%%%%%%%%%%%%%%%%%%%%%%%%%%%%%%%%%%%%%%%%%%%%%%%%%%%%%%%%%

Let \(M\) and \(N\) be connected surfaces.

Remove an open disk from each and identify the resulting boundary circles.

The resulting surface is the connected sum

\[
\boxed{
M\#N.
}
\]

For example,

\[
T^2\#T^2
\]

is the orientable surface of genus \(2\).

%%%%%%%%%%%%%%%%%%%%%%%%%%%%%%%%%%%%%%%%%%%%%%%%%%%%%%%%%%%%
\subsection{Orientable Surfaces of Genus \texorpdfstring{\(g\)}{g}}
\label{subsec:orientable-surfaces-genus}
%%%%%%%%%%%%%%%%%%%%%%%%%%%%%%%%%%%%%%%%%%%%%%%%%%%%%%%%%%%%

Every closed connected orientable surface is homeomorphic to a connected sum
of \(g\) tori:

\[
\boxed{
\Sigma_g
=
\underbrace{
T^2\#\cdots\#T^2
}_{g\text{ copies}}.
}
\]

The integer

\[
g
\]

is called the genus.

For

\[
g=0,
\]

we obtain

\[
\Sigma_0=S^2.
\]

For

\[
g=1,
\]

we obtain

\[
\Sigma_1=T^2.
\]

%%%%%%%%%%%%%%%%%%%%%%%%%%%%%%%%%%%%%%%%%%%%%%%%%%%%%%%%%%%%
\subsection{Euler Characteristic of Orientable Surfaces}
\label{subsec:euler-characteristic-orientable-surfaces}
%%%%%%%%%%%%%%%%%%%%%%%%%%%%%%%%%%%%%%%%%%%%%%%%%%%%%%%%%%%%

For the closed orientable surface of genus \(g\),

\[
\boxed{
\chi(\Sigma_g)=2-2g.
}
\]

Thus

\[
\chi(S^2)=2,
\]

\[
\chi(T^2)=0,
\]

and

\[
\chi(\Sigma_2)=-2.
\]

%%%%%%%%%%%%%%%%%%%%%%%%%%%%%%%%%%%%%%%%%%%%%%%%%%%%%%%%%%%%
\subsection{Worked Example: Determining Genus}
\label{subsec:worked-surface-genus}
%%%%%%%%%%%%%%%%%%%%%%%%%%%%%%%%%%%%%%%%%%%%%%%%%%%%%%%%%%%%

\begin{workedexamplebox}{Genus from Euler Characteristic}

Suppose a closed connected orientable surface has

\[
\chi=-6.
\]

For an orientable surface,

\[
\chi=2-2g.
\]

Therefore

\[
-6=2-2g.
\]

Subtracting \(2\),

\[
-8=-2g.
\]

Hence

\[
g=4.
\]

\begin{answerbox}
\[
\boxed{
g=4.
}
\]
\end{answerbox}
\end{workedexamplebox}

%%%%%%%%%%%%%%%%%%%%%%%%%%%%%%%%%%%%%%%%%%%%%%%%%%%%%%%%%%%%
\subsection{Nonorientable Surfaces}
\label{subsec:nonorientable-surfaces}
%%%%%%%%%%%%%%%%%%%%%%%%%%%%%%%%%%%%%%%%%%%%%%%%%%%%%%%%%%%%

Every closed connected nonorientable surface is homeomorphic to a connected
sum of projective planes:

\[
\boxed{
N_k
=
\underbrace{
\mathbb{RP}^2\#\cdots\#\mathbb{RP}^2
}_{k\text{ copies}}.
}
\]

The integer \(k\) is called the nonorientable genus or crosscap number.

%%%%%%%%%%%%%%%%%%%%%%%%%%%%%%%%%%%%%%%%%%%%%%%%%%%%%%%%%%%%
\subsection{Euler Characteristic of Nonorientable Surfaces}
\label{subsec:euler-characteristic-nonorientable}
%%%%%%%%%%%%%%%%%%%%%%%%%%%%%%%%%%%%%%%%%%%%%%%%%%%%%%%%%%%%

For

\[
N_k,
\]

the Euler characteristic is

\[
\boxed{
\chi(N_k)=2-k.
}
\]

For example,

\[
\chi(\mathbb{RP}^2)=1
\]

and

\[
\chi(\text{Klein bottle})=0.
\]

Indeed, the Klein bottle is homeomorphic to

\[
\boxed{
\mathbb{RP}^2\#\mathbb{RP}^2.
}
\]

%%%%%%%%%%%%%%%%%%%%%%%%%%%%%%%%%%%%%%%%%%%%%%%%%%%%%%%%%%%%
\subsection{Classification of Closed Surfaces}
\label{subsec:classification-closed-surfaces}
%%%%%%%%%%%%%%%%%%%%%%%%%%%%%%%%%%%%%%%%%%%%%%%%%%%%%%%%%%%%

\begin{theorem}[Classification Theorem for Closed Surfaces]
Every compact connected surface without boundary is homeomorphic to exactly
one of the following:

\[
\boxed{
S^2,
}
\]

an orientable surface

\[
\boxed{
\Sigma_g
=
\#^g T^2,
\qquad
g\geq1,
}
\]

or a nonorientable surface

\[
\boxed{
N_k
=
\#^k\mathbb{RP}^2,
\qquad
k\geq1.
}
\]

The orientability and Euler characteristic determine the homeomorphism type.
\end{theorem}

This is one of the most complete classification results in topology.

%%%%%%%%%%%%%%%%%%%%%%%%%%%%%%%%%%%%%%%%%%%%%%%%%%%%%%%%%%%%
\subsection{Surfaces with Boundary}
\label{subsec:surfaces-with-boundary}
%%%%%%%%%%%%%%%%%%%%%%%%%%%%%%%%%%%%%%%%%%%%%%%%%%%%%%%%%%%%

Compact surfaces may also possess boundary components.

For an orientable surface of genus \(g\) with \(b\) boundary components,

\[
\boxed{
\chi=2-2g-b.
}
\]

For a nonorientable surface of crosscap number \(k\) with \(b\) boundary
components,

\[
\boxed{
\chi=2-k-b.
}
\]

%%%%%%%%%%%%%%%%%%%%%%%%%%%%%%%%%%%%%%%%%%%%%%%%%%%%%%%%%%%%
\subsection{Worked Example: Surface with Boundary}
\label{subsec:worked-surface-boundary}
%%%%%%%%%%%%%%%%%%%%%%%%%%%%%%%%%%%%%%%%%%%%%%%%%%%%%%%%%%%%

\begin{workedexamplebox}{Euler Characteristic}

Consider an orientable surface with

\[
g=2
\]

and

\[
b=3
\]

boundary components.

Then

\[
\chi=2-2g-b.
\]

Therefore

\[
\chi
=
2-2(2)-3
=
2-4-3
=
-5.
\]

\begin{answerbox}
\[
\boxed{
\chi=-5.
}
\]
\end{answerbox}
\end{workedexamplebox}

%%%%%%%%%%%%%%%%%%%%%%%%%%%%%%%%%%%%%%%%%%%%%%%%%%%%%%%%%%%%
\subsection{Fundamental Groups of Surfaces}
\label{subsec:fundamental-groups-surfaces}
%%%%%%%%%%%%%%%%%%%%%%%%%%%%%%%%%%%%%%%%%%%%%%%%%%%%%%%%%%%%

The fundamental group provides important algebraic information about a
surface.

For the sphere,

\[
\boxed{
\pi_1(S^2)=\{e\}.
}
\]

For the torus,

\[
\boxed{
\pi_1(T^2)
\cong
\Z\times\Z.
}
\]

For the projective plane,

\[
\boxed{
\pi_1(\mathbb{RP}^2)
\cong
\Z/2\Z.
}
\]

%%%%%%%%%%%%%%%%%%%%%%%%%%%%%%%%%%%%%%%%%%%%%%%%%%%%%%%%%%%%
\subsection{Fundamental Group of an Orientable Surface}
\label{subsec:fundamental-group-orientable-surface}
%%%%%%%%%%%%%%%%%%%%%%%%%%%%%%%%%%%%%%%%%%%%%%%%%%%%%%%%%%%%

For a closed orientable surface of genus \(g\),

\[
\boxed{
\pi_1(\Sigma_g)
=
\left\langle
a_1,b_1,\ldots,a_g,b_g
\;\middle|\;
[a_1,b_1]\cdots[a_g,b_g]=1
\right\rangle.
}
\]

Here

\[
[a,b]
=
aba^{-1}b^{-1}
\]

is the commutator of \(a\) and \(b\).

For \(g=1\),

\[
[a_1,b_1]=1,
\]

so

\[
\pi_1(T^2)
\cong
\Z^2.
\]

%%%%%%%%%%%%%%%%%%%%%%%%%%%%%%%%%%%%%%%%%%%%%%%%%%%%%%%%%%%%
\subsection{Curves on Surfaces}
\label{subsec:curves-on-surfaces}
%%%%%%%%%%%%%%%%%%%%%%%%%%%%%%%%%%%%%%%%%%%%%%%%%%%%%%%%%%%%

Simple closed curves are central objects in surface topology.

A simple closed curve is an embedding

\[
\boxed{
S^1\hookrightarrow\Sigma.
}
\]

Such a curve may be

\[
\boxed{
\text{contractible},
\qquad
\text{essential},
\qquad
\text{separating},
\qquad
\text{nonseparating}.
}
\]

%%%%%%%%%%%%%%%%%%%%%%%%%%%%%%%%%%%%%%%%%%%%%%%%%%%%%%%%%%%%
\subsection{Separating Curves}
\label{subsec:separating-curves}
%%%%%%%%%%%%%%%%%%%%%%%%%%%%%%%%%%%%%%%%%%%%%%%%%%%%%%%%%%%%

A simple closed curve

\[
\gamma\subset\Sigma
\]

is separating if

\[
\boxed{
\Sigma\setminus\gamma
}
\]

is disconnected.

Otherwise it is nonseparating.

On a torus, an essential simple closed curve is nonseparating.

On a genus-\(2\) surface, both separating and nonseparating essential curves
occur.

%%%%%%%%%%%%%%%%%%%%%%%%%%%%%%%%%%%%%%%%%%%%%%%%%%%%%%%%%%%%
\subsection{Intersection of Curves}
\label{subsec:intersection-curves-surfaces}
%%%%%%%%%%%%%%%%%%%%%%%%%%%%%%%%%%%%%%%%%%%%%%%%%%%%%%%%%%%%

Two curves on a surface may intersect.

After deforming them within their isotopy classes, one can attempt to
minimize the number of intersections.

This leads to the geometric intersection number

\[
\boxed{
i(\alpha,\beta).
}
\]

It is the minimum number of intersection points among representatives of the
two isotopy classes.

%%%%%%%%%%%%%%%%%%%%%%%%%%%%%%%%%%%%%%%%%%%%%%%%%%%%%%%%%%%%
\subsection{Homeomorphisms of Surfaces}
\label{subsec:homeomorphisms-surfaces}
%%%%%%%%%%%%%%%%%%%%%%%%%%%%%%%%%%%%%%%%%%%%%%%%%%%%%%%%%%%%

A homeomorphism

\[
f:\Sigma\to\Sigma
\]

may twist, rotate, or otherwise rearrange a surface.

Instead of distinguishing homeomorphisms that differ only by a continuous
deformation, one studies isotopy classes of homeomorphisms.

This produces the mapping class group.

%%%%%%%%%%%%%%%%%%%%%%%%%%%%%%%%%%%%%%%%%%%%%%%%%%%%%%%%%%%%
\subsection{Mapping Class Group}
\label{subsec:mapping-class-group}
%%%%%%%%%%%%%%%%%%%%%%%%%%%%%%%%%%%%%%%%%%%%%%%%%%%%%%%%%%%%

\begin{definitionbox}{Mapping Class Group}
For an oriented surface \(\Sigma\), the orientation-preserving mapping class
group is

\[
\boxed{
\operatorname{Mod}(\Sigma)
=
\operatorname{Homeo}^{+}(\Sigma)
/
\operatorname{Homeo}_{0}(\Sigma),
}
\]

where \(\operatorname{Homeo}_{0}(\Sigma)\) denotes the subgroup of
homeomorphisms isotopic to the identity.
\end{definitionbox}

Thus the mapping class group records topologically distinct symmetries of
the surface.

%%%%%%%%%%%%%%%%%%%%%%%%%%%%%%%%%%%%%%%%%%%%%%%%%%%%%%%%%%%%
\subsection{Dehn Twists}
\label{subsec:dehn-twists}
%%%%%%%%%%%%%%%%%%%%%%%%%%%%%%%%%%%%%%%%%%%%%%%%%%%%%%%%%%%%

A Dehn twist is a fundamental homeomorphism of a surface.

Choose a simple closed curve

\[
\gamma\subset\Sigma.
\]

Cut along \(\gamma\), rotate one side through one full twist, and glue the
surface back together.

The resulting mapping class is denoted by notation such as

\[
\boxed{
T_\gamma.
}
\]

Dehn twists generate the mapping class group of a compact orientable
surface.

%%%%%%%%%%%%%%%%%%%%%%%%%%%%%%%%%%%%%%%%%%%%%%%%%%%%%%%%%%%%
\subsection{Dimension Three}
\label{subsec:dimension-three-topology}
%%%%%%%%%%%%%%%%%%%%%%%%%%%%%%%%%%%%%%%%%%%%%%%%%%%%%%%%%%%%

A \(3\)-manifold is a space locally modeled on

\[
\R^3.
\]

Examples include

\[
\boxed{
S^3,
\qquad
S^2\times S^1,
\qquad
T^3,
}
\]

as well as knot complements and lens spaces.

Three-dimensional topology combines

\[
\boxed{
\text{topology},
\quad
\text{geometry},
\quad
\text{group theory},
\quad
\text{knot theory}.
}
\]

%%%%%%%%%%%%%%%%%%%%%%%%%%%%%%%%%%%%%%%%%%%%%%%%%%%%%%%%%%%%
\subsection{The Three-Sphere}
\label{subsec:three-sphere-low-dimensional}
%%%%%%%%%%%%%%%%%%%%%%%%%%%%%%%%%%%%%%%%%%%%%%%%%%%%%%%%%%%%

The \(3\)-sphere is

\[
\boxed{
S^3
=
\left\{
(x_1,x_2,x_3,x_4)\in\R^4:
x_1^2+x_2^2+x_3^2+x_4^2=1
\right\}.
}
\]

It is the three-dimensional analogue of the ordinary sphere.

It is compact, connected, orientable, and simply connected.

%%%%%%%%%%%%%%%%%%%%%%%%%%%%%%%%%%%%%%%%%%%%%%%%%%%%%%%%%%%%
\subsection{Lens Spaces}
\label{subsec:lens-spaces}
%%%%%%%%%%%%%%%%%%%%%%%%%%%%%%%%%%%%%%%%%%%%%%%%%%%%%%%%%%%%

Lens spaces form an important family of closed \(3\)-manifolds.

They are commonly denoted

\[
\boxed{
L(p,q).
}
\]

They can be constructed through quotients of \(S^3\), gluing solid tori, or
Dehn surgery on the unknot.

For

\[
p>0,
\]

one has

\[
\boxed{
\pi_1(L(p,q))
\cong
\Z/p\Z.
}
\]

Different lens spaces can have the same fundamental group, illustrating that
fundamental groups alone do not completely classify \(3\)-manifolds.

%%%%%%%%%%%%%%%%%%%%%%%%%%%%%%%%%%%%%%%%%%%%%%%%%%%%%%%%%%%%
\subsection{Heegaard Splittings}
\label{subsec:heegaard-splittings}
%%%%%%%%%%%%%%%%%%%%%%%%%%%%%%%%%%%%%%%%%%%%%%%%%%%%%%%%%%%%

A powerful way to describe a closed orientable \(3\)-manifold is to split it
into two handlebodies.

\begin{definitionbox}{Heegaard Splitting}
A Heegaard splitting of a closed orientable \(3\)-manifold \(M\) is a
decomposition

\[
\boxed{
M=H_1\cup_{\Sigma_g}H_2,
}
\]

where \(H_1\) and \(H_2\) are genus-\(g\) handlebodies and

\[
\partial H_1
=
\partial H_2
=
\Sigma_g.
\]
\end{definitionbox}

The surface

\[
\Sigma_g
\]

is called a Heegaard surface.

%%%%%%%%%%%%%%%%%%%%%%%%%%%%%%%%%%%%%%%%%%%%%%%%%%%%%%%%%%%%
\subsection{Heegaard Genus}
\label{subsec:heegaard-genus}
%%%%%%%%%%%%%%%%%%%%%%%%%%%%%%%%%%%%%%%%%%%%%%%%%%%%%%%%%%%%

The Heegaard genus of a closed orientable \(3\)-manifold \(M\) is the
smallest genus of any Heegaard splitting of \(M\).

It is denoted

\[
\boxed{
g(M).
}
\]

For the \(3\)-sphere,

\[
\boxed{
g(S^3)=0.
}
\]

Lens spaces have Heegaard genus at most \(1\), and nontrivial lens spaces
have genus \(1\).

%%%%%%%%%%%%%%%%%%%%%%%%%%%%%%%%%%%%%%%%%%%%%%%%%%%%%%%%%%%%
\subsection{Handlebodies}
\label{subsec:handlebodies-low-dimensional}
%%%%%%%%%%%%%%%%%%%%%%%%%%%%%%%%%%%%%%%%%%%%%%%%%%%%%%%%%%%%

A genus-\(g\) handlebody can be obtained from a \(3\)-ball by attaching
\(g\) one-handles.

Its boundary is

\[
\boxed{
\partial H_g=\Sigma_g.
}
\]

For \(g=0\),

\[
H_0=D^3.
\]

For \(g=1\),

\[
H_1
\]

is a solid torus.

%%%%%%%%%%%%%%%%%%%%%%%%%%%%%%%%%%%%%%%%%%%%%%%%%%%%%%%%%%%%
\subsection{Every Closed Orientable Three-Manifold Has a Heegaard Splitting}
\label{subsec:existence-heegaard}
%%%%%%%%%%%%%%%%%%%%%%%%%%%%%%%%%%%%%%%%%%%%%%%%%%%%%%%%%%%%

A fundamental theorem states that every closed orientable \(3\)-manifold
admits a Heegaard splitting.

Thus complicated \(3\)-manifolds can be encoded by

\[
\boxed{
\text{surface}
+
\text{gluing homeomorphism}.
}
\]

This makes mapping class groups directly relevant to \(3\)-manifold
topology.

%%%%%%%%%%%%%%%%%%%%%%%%%%%%%%%%%%%%%%%%%%%%%%%%%%%%%%%%%%%%
\subsection{Triangulations of Three-Manifolds}
\label{subsec:triangulations-three-manifolds}
%%%%%%%%%%%%%%%%%%%%%%%%%%%%%%%%%%%%%%%%%%%%%%%%%%%%%%%%%%%%

A \(3\)-manifold can often be represented by gluing tetrahedra along their
faces.

Thus a manifold may be encoded combinatorially by

\[
\boxed{
\text{tetrahedra}
+
\text{face identifications}.
}
\]

Triangulations make computational approaches to \(3\)-manifold topology
possible.

%%%%%%%%%%%%%%%%%%%%%%%%%%%%%%%%%%%%%%%%%%%%%%%%%%%%%%%%%%%%
\subsection{Prime Three-Manifolds}
\label{subsec:prime-three-manifolds}
%%%%%%%%%%%%%%%%%%%%%%%%%%%%%%%%%%%%%%%%%%%%%%%%%%%%%%%%%%%%

A connected orientable \(3\)-manifold \(M\) is prime if

\[
M=M_1\#M_2
\]

implies that one of the summands is homeomorphic to \(S^3\).

Prime manifolds play a role analogous to prime knots and prime integers.

%%%%%%%%%%%%%%%%%%%%%%%%%%%%%%%%%%%%%%%%%%%%%%%%%%%%%%%%%%%%
\subsection{Prime Decomposition}
\label{subsec:prime-decomposition-three-manifolds}
%%%%%%%%%%%%%%%%%%%%%%%%%%%%%%%%%%%%%%%%%%%%%%%%%%%%%%%%%%%%

Closed orientable \(3\)-manifolds admit a decomposition into prime pieces.

Schematically,

\[
\boxed{
M
=
M_1\#M_2\#\cdots\#M_k.
}
\]

Under the standard hypotheses, this decomposition is essentially unique up
to order and homeomorphism of the prime factors.

This reduces the study of general \(3\)-manifolds to prime manifolds.

%%%%%%%%%%%%%%%%%%%%%%%%%%%%%%%%%%%%%%%%%%%%%%%%%%%%%%%%%%%%
\subsection{Irreducible Three-Manifolds}
\label{subsec:irreducible-three-manifolds}
%%%%%%%%%%%%%%%%%%%%%%%%%%%%%%%%%%%%%%%%%%%%%%%%%%%%%%%%%%%%

A \(3\)-manifold is irreducible if every embedded \(2\)-sphere bounds a
\(3\)-ball.

Thus

\[
\boxed{
S^2\hookrightarrow M
}
\]

must be topologically trivial in the appropriate sense.

Most prime orientable \(3\)-manifolds are irreducible, with

\[
S^2\times S^1
\]

providing the standard exceptional prime example.

%%%%%%%%%%%%%%%%%%%%%%%%%%%%%%%%%%%%%%%%%%%%%%%%%%%%%%%%%%%%
\subsection{Incompressible Surfaces}
\label{subsec:incompressible-surfaces}
%%%%%%%%%%%%%%%%%%%%%%%%%%%%%%%%%%%%%%%%%%%%%%%%%%%%%%%%%%%%

Embedded surfaces provide essential tools for decomposing \(3\)-manifolds.

Roughly, an embedded surface

\[
F\subset M
\]

is incompressible if it cannot be simplified by compression along a disk.

Under standard hypotheses, incompressibility can be characterized through
injectivity of

\[
\boxed{
\pi_1(F)\longrightarrow\pi_1(M).
}
\]

These surfaces reveal internal structure of the ambient manifold.

%%%%%%%%%%%%%%%%%%%%%%%%%%%%%%%%%%%%%%%%%%%%%%%%%%%%%%%%%%%%
\subsection{Haken Manifolds}
\label{subsec:haken-manifolds}
%%%%%%%%%%%%%%%%%%%%%%%%%%%%%%%%%%%%%%%%%%%%%%%%%%%%%%%%%%%%

A Haken manifold is, roughly, a sufficiently well-behaved irreducible
\(3\)-manifold containing a properly embedded two-sided incompressible
surface.

Such surfaces allow the manifold to be cut into simpler pieces.

This idea was foundational in the development of algorithmic and structural
\(3\)-manifold topology.

%%%%%%%%%%%%%%%%%%%%%%%%%%%%%%%%%%%%%%%%%%%%%%%%%%%%%%%%%%%%
\subsection{JSJ Decomposition}
\label{subsec:jsj-decomposition}
%%%%%%%%%%%%%%%%%%%%%%%%%%%%%%%%%%%%%%%%%%%%%%%%%%%%%%%%%%%%

Certain irreducible \(3\)-manifolds can be cut along essential embedded tori
into simpler pieces.

This is the JSJ decomposition.

Schematically,

\[
\boxed{
M
\longrightarrow
\text{cut along essential tori}
\longrightarrow
\text{geometric pieces}.
}
\]

The initials JSJ refer to Jaco--Shalen and Johannson.

The decomposition plays a central role in understanding the structure of
\(3\)-manifolds.

%%%%%%%%%%%%%%%%%%%%%%%%%%%%%%%%%%%%%%%%%%%%%%%%%%%%%%%%%%%%
\subsection{Thurston's Geometrization Program}
\label{subsec:thurston-geometrization}
%%%%%%%%%%%%%%%%%%%%%%%%%%%%%%%%%%%%%%%%%%%%%%%%%%%%%%%%%%%%

Thurston proposed that suitable pieces of \(3\)-manifolds should admit
canonical geometric structures.

In this viewpoint, topology is understood by decomposing a manifold and
assigning geometry to its pieces.

The possible model geometries form a finite list.

%%%%%%%%%%%%%%%%%%%%%%%%%%%%%%%%%%%%%%%%%%%%%%%%%%%%%%%%%%%%
\subsection{The Eight Thurston Geometries}
\label{subsec:eight-thurston-geometries}
%%%%%%%%%%%%%%%%%%%%%%%%%%%%%%%%%%%%%%%%%%%%%%%%%%%%%%%%%%%%

The eight model geometries are

\[
\boxed{
S^3,
\qquad
\mathbb{E}^3,
\qquad
\mathbb{H}^3,
}
\]

\[
\boxed{
S^2\times\R,
\qquad
\mathbb{H}^2\times\R,
}
\]

\[
\boxed{
\widetilde{\mathrm{SL}_2(\R)},
\qquad
\mathrm{Nil},
\qquad
\mathrm{Sol}.
}
\]

Here

\[
\mathbb{E}^3
\]

denotes Euclidean \(3\)-space and

\[
\mathbb{H}^3
\]

denotes hyperbolic \(3\)-space.

%%%%%%%%%%%%%%%%%%%%%%%%%%%%%%%%%%%%%%%%%%%%%%%%%%%%%%%%%%%%
\subsection{Hyperbolic Three-Manifolds}
\label{subsec:hyperbolic-three-manifolds}
%%%%%%%%%%%%%%%%%%%%%%%%%%%%%%%%%%%%%%%%%%%%%%%%%%%%%%%%%%%%

A hyperbolic \(3\)-manifold locally carries the geometry of

\[
\mathbb{H}^3.
\]

Many knot complements are hyperbolic.

The figure-eight knot complement is a standard example.

For finite-volume hyperbolic \(3\)-manifolds, Mostow--Prasad rigidity makes
geometric quantities such as volume topological invariants.

%%%%%%%%%%%%%%%%%%%%%%%%%%%%%%%%%%%%%%%%%%%%%%%%%%%%%%%%%%%%
\subsection{Geometrization Theorem}
\label{subsec:geometrization-theorem}
%%%%%%%%%%%%%%%%%%%%%%%%%%%%%%%%%%%%%%%%%%%%%%%%%%%%%%%%%%%%

The geometrization theorem, proved through the work initiated by Richard
Hamilton and completed by Grigori Perelman using Ricci flow, establishes the
geometric decomposition envisioned by Thurston.

Conceptually,

\[
\boxed{
\text{topological decomposition}
+
\text{geometric structures}
}
\]

provides a powerful framework for understanding closed \(3\)-manifolds.

%%%%%%%%%%%%%%%%%%%%%%%%%%%%%%%%%%%%%%%%%%%%%%%%%%%%%%%%%%%%
\subsection{The Poincar\'e Conjecture}
\label{subsec:poincare-conjecture}
%%%%%%%%%%%%%%%%%%%%%%%%%%%%%%%%%%%%%%%%%%%%%%%%%%%%%%%%%%%%

The Poincar\'e conjecture asked whether every closed simply connected
\(3\)-manifold is homeomorphic to the \(3\)-sphere.

In modern form:

\begin{theorem}[Poincar\'e Conjecture]
If \(M\) is a closed \(3\)-manifold and

\[
\pi_1(M)=\{e\},
\]

then

\[
\boxed{
M\cong S^3.
}
\]
\end{theorem}

The theorem follows from Perelman's proof of geometrization.

%%%%%%%%%%%%%%%%%%%%%%%%%%%%%%%%%%%%%%%%%%%%%%%%%%%%%%%%%%%%
\subsection{Ricci Flow}
\label{subsec:ricci-flow-low-dimensional}
%%%%%%%%%%%%%%%%%%%%%%%%%%%%%%%%%%%%%%%%%%%%%%%%%%%%%%%%%%%%

Ricci flow evolves a Riemannian metric according to

\[
\boxed{
\frac{\partial g}{\partial t}
=
-2\operatorname{Ric}(g).
}
\]

Here

\[
\operatorname{Ric}
\]

denotes the Ricci curvature tensor.

The equation tends to smooth geometric irregularities while developing
singularities that must be carefully analyzed.

Ricci flow with surgery was central to the proof of geometrization.

%%%%%%%%%%%%%%%%%%%%%%%%%%%%%%%%%%%%%%%%%%%%%%%%%%%%%%%%%%%%
\subsection{Dehn Surgery and Three-Manifolds}
\label{subsec:dehn-surgery-three-manifolds}
%%%%%%%%%%%%%%%%%%%%%%%%%%%%%%%%%%%%%%%%%%%%%%%%%%%%%%%%%%%%

Knot theory provides another method for constructing \(3\)-manifolds.

Begin with a link

\[
L\subset S^3.
\]

Remove tubular neighborhoods of its components and glue solid tori back
using specified boundary identifications.

The result is a new closed \(3\)-manifold.

This is surgery on a framed link.

%%%%%%%%%%%%%%%%%%%%%%%%%%%%%%%%%%%%%%%%%%%%%%%%%%%%%%%%%%%%
\subsection{Surgery Presentations}
\label{subsec:surgery-presentations}
%%%%%%%%%%%%%%%%%%%%%%%%%%%%%%%%%%%%%%%%%%%%%%%%%%%%%%%%%%%%

A framed link diagram can encode a \(3\)-manifold.

Thus

\[
\boxed{
\text{framed link diagram}
\longrightarrow
\text{closed }3\text{-manifold}.
}
\]

This creates a powerful bridge between knot diagrams and manifold topology.

%%%%%%%%%%%%%%%%%%%%%%%%%%%%%%%%%%%%%%%%%%%%%%%%%%%%%%%%%%%%
\subsection{Lickorish--Wallace Theorem}
\label{subsec:lickorish-wallace}
%%%%%%%%%%%%%%%%%%%%%%%%%%%%%%%%%%%%%%%%%%%%%%%%%%%%%%%%%%%%

\begin{theorem}[Lickorish--Wallace Theorem]
Every closed connected orientable \(3\)-manifold can be obtained by Dehn
surgery on a link in \(S^3\).
\end{theorem}

Therefore link diagrams can, in principle, encode every closed orientable
\(3\)-manifold.

%%%%%%%%%%%%%%%%%%%%%%%%%%%%%%%%%%%%%%%%%%%%%%%%%%%%%%%%%%%%
\subsection{Dimension Four}
\label{subsec:dimension-four-topology}
%%%%%%%%%%%%%%%%%%%%%%%%%%%%%%%%%%%%%%%%%%%%%%%%%%%%%%%%%%%%

Dimension four is one of the most unusual dimensions in mathematics.

A \(4\)-manifold is locally modeled on

\[
\R^4.
\]

However, the topological and smooth classifications of \(4\)-manifolds
behave very differently.

This produces phenomena that do not occur in the same way in other
dimensions.

%%%%%%%%%%%%%%%%%%%%%%%%%%%%%%%%%%%%%%%%%%%%%%%%%%%%%%%%%%%%
\subsection{Topological Versus Smooth Manifolds}
\label{subsec:topological-versus-smooth-four-manifolds}
%%%%%%%%%%%%%%%%%%%%%%%%%%%%%%%%%%%%%%%%%%%%%%%%%%%%%%%%%%%%

A topological manifold is defined using homeomorphisms.

A smooth manifold additionally possesses coordinate charts with smooth
transition maps.

In many dimensions, the relationship between these structures is relatively
well controlled.

In dimension four, it is exceptionally complicated.

One topological \(4\)-manifold may support multiple inequivalent smooth
structures.

%%%%%%%%%%%%%%%%%%%%%%%%%%%%%%%%%%%%%%%%%%%%%%%%%%%%%%%%%%%%
\subsection{Exotic Smooth Structures}
\label{subsec:exotic-smooth-structures}
%%%%%%%%%%%%%%%%%%%%%%%%%%%%%%%%%%%%%%%%%%%%%%%%%%%%%%%%%%%%

Two smooth manifolds may be homeomorphic but not diffeomorphic.

When this occurs, they have the same underlying topology but different
smooth structures.

Such structures are called exotic.

Dimension four is particularly remarkable because

\[
\boxed{
\R^4
}
\]

admits exotic smooth structures.

%%%%%%%%%%%%%%%%%%%%%%%%%%%%%%%%%%%%%%%%%%%%%%%%%%%%%%%%%%%%
\subsection{Exotic \texorpdfstring{\(\R^4\)}{R4}}
\label{subsec:exotic-r4}
%%%%%%%%%%%%%%%%%%%%%%%%%%%%%%%%%%%%%%%%%%%%%%%%%%%%%%%%%%%%

There exist smooth \(4\)-manifolds that are homeomorphic to ordinary

\[
\R^4
\]

but not diffeomorphic to it.

These are called exotic \(\R^4\)'s.

Thus the same topological space can support genuinely different notions of
smoothness.

This phenomenon is unique to dimension four among Euclidean spaces.

%%%%%%%%%%%%%%%%%%%%%%%%%%%%%%%%%%%%%%%%%%%%%%%%%%%%%%%%%%%%
\subsection{Surfaces Inside Four-Manifolds}
\label{subsec:surfaces-inside-four-manifolds}
%%%%%%%%%%%%%%%%%%%%%%%%%%%%%%%%%%%%%%%%%%%%%%%%%%%%%%%%%%%%

A closed oriented surface can be embedded inside an oriented \(4\)-manifold:

\[
\boxed{
\Sigma\hookrightarrow M^4.
}
\]

Because

\[
2+2=4,
\]

two such surfaces generically intersect in isolated points.

These intersections can be counted algebraically.

This leads to the intersection form.

%%%%%%%%%%%%%%%%%%%%%%%%%%%%%%%%%%%%%%%%%%%%%%%%%%%%%%%%%%%%
\subsection{Intersection Form}
\label{subsec:intersection-form}
%%%%%%%%%%%%%%%%%%%%%%%%%%%%%%%%%%%%%%%%%%%%%%%%%%%%%%%%%%%%

For a closed oriented \(4\)-manifold \(M\), the intersection form is a
bilinear pairing

\[
\boxed{
Q_M:
H_2(M;\Z)/\operatorname{Tor}
\times
H_2(M;\Z)/\operatorname{Tor}
\to
\Z.
}
\]

It records algebraic intersections between representatives of
two-dimensional homology classes.

Here

\[
\operatorname{Tor}
\]

denotes the torsion subgroup.

%%%%%%%%%%%%%%%%%%%%%%%%%%%%%%%%%%%%%%%%%%%%%%%%%%%%%%%%%%%%
\subsection{Matrix of an Intersection Form}
\label{subsec:intersection-form-matrix}
%%%%%%%%%%%%%%%%%%%%%%%%%%%%%%%%%%%%%%%%%%%%%%%%%%%%%%%%%%%%

Choose a basis

\[
e_1,\ldots,e_n
\]

for the free part of

\[
H_2(M;\Z).
\]

Then the intersection form can be represented by a symmetric integer matrix

\[
\boxed{
Q=(Q_{ij}),
}
\]

where

\[
Q_{ij}
=
Q_M(e_i,e_j).
\]

Thus the topology of a \(4\)-manifold produces an object from integral
linear algebra.

%%%%%%%%%%%%%%%%%%%%%%%%%%%%%%%%%%%%%%%%%%%%%%%%%%%%%%%%%%%%
\subsection{Examples of Intersection Forms}
\label{subsec:examples-intersection-forms}
%%%%%%%%%%%%%%%%%%%%%%%%%%%%%%%%%%%%%%%%%%%%%%%%%%%%%%%%%%%%

For the complex projective plane,

\[
\mathbb{CP}^2,
\]

the intersection form is represented by

\[
\boxed{
[1].
}
\]

For the oppositely oriented complex projective plane,

\[
\overline{\mathbb{CP}}^{\,2},
\]

the form is

\[
\boxed{
[-1].
}
\]

For

\[
S^2\times S^2,
\]

one obtains the hyperbolic form

\[
\boxed{
H
=
\begin{pmatrix}
0&1\\
1&0
\end{pmatrix}.
}
\]

%%%%%%%%%%%%%%%%%%%%%%%%%%%%%%%%%%%%%%%%%%%%%%%%%%%%%%%%%%%%
\subsection{Signature}
\label{subsec:four-manifold-signature}
%%%%%%%%%%%%%%%%%%%%%%%%%%%%%%%%%%%%%%%%%%%%%%%%%%%%%%%%%%%%

After extending an intersection form to the real numbers, its matrix has
positive and negative eigenvalues.

If

\[
b_2^+
\]

is the number of positive eigenvalues and

\[
b_2^-
\]

the number of negative eigenvalues, the signature is

\[
\boxed{
\sigma(M)
=
b_2^+-b_2^-.
}
\]

The lowercase Greek letter

\[
\sigma
\]

is read ``sigma.''

%%%%%%%%%%%%%%%%%%%%%%%%%%%%%%%%%%%%%%%%%%%%%%%%%%%%%%%%%%%%
\subsection{Definite and Indefinite Forms}
\label{subsec:definite-indefinite-forms}
%%%%%%%%%%%%%%%%%%%%%%%%%%%%%%%%%%%%%%%%%%%%%%%%%%%%%%%%%%%%

An intersection form is positive definite if

\[
Q(x,x)>0
\]

for every nonzero \(x\).

It is negative definite if

\[
Q(x,x)<0
\]

for every nonzero \(x\).

Otherwise, when both positive and negative directions occur, the form is
indefinite.

These distinctions have major consequences in \(4\)-manifold topology.

%%%%%%%%%%%%%%%%%%%%%%%%%%%%%%%%%%%%%%%%%%%%%%%%%%%%%%%%%%%%
\subsection{Freedman's Theorem}
\label{subsec:freedman-theorem}
%%%%%%%%%%%%%%%%%%%%%%%%%%%%%%%%%%%%%%%%%%%%%%%%%%%%%%%%%%%%

Freedman's work established a deep classification theory for simply
connected topological \(4\)-manifolds.

Very roughly, intersection forms and additional data provide powerful
classification information in the topological category.

This result revealed that topological \(4\)-manifolds can be substantially
more flexible than smooth \(4\)-manifolds.

%%%%%%%%%%%%%%%%%%%%%%%%%%%%%%%%%%%%%%%%%%%%%%%%%%%%%%%%%%%%
\subsection{Donald's Theorem}
\label{subsec:donaldson-theorem}
%%%%%%%%%%%%%%%%%%%%%%%%%%%%%%%%%%%%%%%%%%%%%%%%%%%%%%%%%%%%

Donaldson's work used gauge theory to establish strong restrictions on
intersection forms of smooth \(4\)-manifolds.

One celebrated consequence is that a definite unimodular intersection form
of a smooth closed simply connected \(4\)-manifold must be diagonalizable
over the integers.

Thus

\[
\boxed{
\text{topologically possible}
}
\]

does not always imply

\[
\boxed{
\text{smoothly possible}.
}
\]

%%%%%%%%%%%%%%%%%%%%%%%%%%%%%%%%%%%%%%%%%%%%%%%%%%%%%%%%%%%%
\subsection{The Four-Dimensional Contrast}
\label{subsec:four-dimensional-contrast}
%%%%%%%%%%%%%%%%%%%%%%%%%%%%%%%%%%%%%%%%%%%%%%%%%%%%%%%%%%%%

The combination of Freedman's and Donaldson's theories exposed a dramatic
difference between

\[
\boxed{
\text{topological }4\text{-manifolds}
}
\]

and

\[
\boxed{
\text{smooth }4\text{-manifolds}.
}
\]

This contrast is one of the defining features of modern
low-dimensional topology.

%%%%%%%%%%%%%%%%%%%%%%%%%%%%%%%%%%%%%%%%%%%%%%%%%%%%%%%%%%%%
\subsection{Handles}
\label{subsec:handles-low-dimensional}
%%%%%%%%%%%%%%%%%%%%%%%%%%%%%%%%%%%%%%%%%%%%%%%%%%%%%%%%%%%%

An \(n\)-dimensional \(k\)-handle is

\[
\boxed{
D^k\times D^{n-k}.
}
\]

It is attached along part of its boundary,

\[
\boxed{
S^{k-1}\times D^{n-k}.
}
\]

A manifold can often be built by attaching handles in increasing order of
index.

%%%%%%%%%%%%%%%%%%%%%%%%%%%%%%%%%%%%%%%%%%%%%%%%%%%%%%%%%%%%
\subsection{Handles in Four Dimensions}
\label{subsec:four-dimensional-handles}
%%%%%%%%%%%%%%%%%%%%%%%%%%%%%%%%%%%%%%%%%%%%%%%%%%%%%%%%%%%%

For a \(4\)-manifold, the possible handles are

\[
\boxed{
0\text{-handles},
\quad
1\text{-handles},
\quad
2\text{-handles},
\quad
3\text{-handles},
\quad
4\text{-handles}.
}
\]

A \(2\)-handle has the form

\[
\boxed{
D^2\times D^2
}
\]

and is attached along

\[
\boxed{
S^1\times D^2.
}
\]

Knots and framed links therefore naturally encode \(2\)-handle attachments.

%%%%%%%%%%%%%%%%%%%%%%%%%%%%%%%%%%%%%%%%%%%%%%%%%%%%%%%%%%%%
\subsection{Kirby Diagrams}
\label{subsec:kirby-diagrams}
%%%%%%%%%%%%%%%%%%%%%%%%%%%%%%%%%%%%%%%%%%%%%%%%%%%%%%%%%%%%

Kirby diagrams provide diagrammatic descriptions of smooth \(4\)-manifolds.

Framed knots represent \(2\)-handles.

Dotted circles are commonly used to encode \(1\)-handles.

Thus a four-dimensional manifold can be studied through a two-dimensional
link diagram.

Schematically,

\[
\boxed{
\text{link diagram}
\longrightarrow
\text{handle decomposition}
\longrightarrow
4\text{-manifold}.
}
\]

%%%%%%%%%%%%%%%%%%%%%%%%%%%%%%%%%%%%%%%%%%%%%%%%%%%%%%%%%%%%
\subsection{Kirby Calculus}
\label{subsec:kirby-calculus}
%%%%%%%%%%%%%%%%%%%%%%%%%%%%%%%%%%%%%%%%%%%%%%%%%%%%%%%%%%%%

Different Kirby diagrams may describe diffeomorphic \(4\)-manifolds.

Kirby calculus provides moves that transform one handle description into
another without changing the underlying manifold.

This plays a role for \(4\)-manifolds analogous to the role played by

\[
\boxed{
\text{Reidemeister moves}
}
\]

for knots.

%%%%%%%%%%%%%%%%%%%%%%%%%%%%%%%%%%%%%%%%%%%%%%%%%%%%%%%%%%%%
\subsection{Knots and Four-Manifolds}
\label{subsec:knots-four-manifolds}
%%%%%%%%%%%%%%%%%%%%%%%%%%%%%%%%%%%%%%%%%%%%%%%%%%%%%%%%%%%%

A knot

\[
K\subset S^3
\]

can be viewed as the boundary of a surface inside

\[
D^4.
\]

The central question becomes:

\[
\boxed{
\text{What surfaces can }K\text{ bound in }D^4?
}
\]

This connects classical knot theory with four-dimensional topology.

%%%%%%%%%%%%%%%%%%%%%%%%%%%%%%%%%%%%%%%%%%%%%%%%%%%%%%%%%%%%
\subsection{Slice Genus}
\label{subsec:slice-genus}
%%%%%%%%%%%%%%%%%%%%%%%%%%%%%%%%%%%%%%%%%%%%%%%%%%%%%%%%%%%%

\begin{definitionbox}{Slice Genus}
The smooth slice genus of a knot \(K\), commonly denoted

\[
\boxed{
g_4(K),
}
\]

is the minimum genus of a smooth, compact, connected, orientable surface

\[
F\subset D^4
\]

with

\[
\partial F=K.
\]
\end{definitionbox}

A knot is smoothly slice precisely when

\[
\boxed{
g_4(K)=0.
}
\]

%%%%%%%%%%%%%%%%%%%%%%%%%%%%%%%%%%%%%%%%%%%%%%%%%%%%%%%%%%%%
\subsection{Three-Genus Versus Four-Genus}
\label{subsec:three-genus-four-genus}
%%%%%%%%%%%%%%%%%%%%%%%%%%%%%%%%%%%%%%%%%%%%%%%%%%%%%%%%%%%%

The ordinary knot genus

\[
g(K)
\]

uses surfaces in \(S^3\).

The slice genus

\[
g_4(K)
\]

uses surfaces in \(D^4\).

Since a Seifert surface can be pushed slightly into the \(4\)-ball,

\[
\boxed{
g_4(K)\leq g(K).
}
\]

The inequality may be strict.

%%%%%%%%%%%%%%%%%%%%%%%%%%%%%%%%%%%%%%%%%%%%%%%%%%%%%%%%%%%%
\subsection{Worked Example: Trefoil Genera}
\label{subsec:worked-trefoil-genera}
%%%%%%%%%%%%%%%%%%%%%%%%%%%%%%%%%%%%%%%%%%%%%%%%%%%%%%%%%%%%

\begin{workedexamplebox}{Three-Genus and Four-Genus}

For the trefoil knot \(3_1\),

\[
g(3_1)=1.
\]

Its smooth slice genus is also

\[
g_4(3_1)=1.
\]

Therefore the trefoil is not slice.

\begin{answerbox}
\[
\boxed{
g(3_1)=g_4(3_1)=1.
}
\]
\end{answerbox}
\end{workedexamplebox}

%%%%%%%%%%%%%%%%%%%%%%%%%%%%%%%%%%%%%%%%%%%%%%%%%%%%%%%%%%%%
\subsection{Concordance Revisited}
\label{subsec:concordance-low-dimensional}
%%%%%%%%%%%%%%%%%%%%%%%%%%%%%%%%%%%%%%%%%%%%%%%%%%%%%%%%%%%%

Two knots are concordant if they cobound a properly embedded annulus in

\[
S^3\times[0,1].
\]

Concordance asks whether knots become equivalent after allowing movement
through a four-dimensional cobordism.

Thus

\[
\boxed{
\text{knot theory in }S^3
}
\]

naturally extends to

\[
\boxed{
\text{surface theory in }D^4.
}
\]

%%%%%%%%%%%%%%%%%%%%%%%%%%%%%%%%%%%%%%%%%%%%%%%%%%%%%%%%%%%%
\subsection{Cobordism}
\label{subsec:cobordism-low-dimensional}
%%%%%%%%%%%%%%%%%%%%%%%%%%%%%%%%%%%%%%%%%%%%%%%%%%%%%%%%%%%%

Cobordism provides a general framework for relating manifolds.

\begin{definitionbox}{Cobordism}
Two closed \(n\)-manifolds \(M_0\) and \(M_1\) are cobordant if there exists
a compact \((n+1)\)-manifold \(W\) whose boundary is

\[
\boxed{
\partial W
=
M_0\sqcup M_1
}
\]

in the unoriented setting, with the appropriate orientation convention in
the oriented setting.
\end{definitionbox}

One may think of \(W\) as a geometric interpolation between its boundary
manifolds.

%%%%%%%%%%%%%%%%%%%%%%%%%%%%%%%%%%%%%%%%%%%%%%%%%%%%%%%%%%%%
\subsection{Morse Theory and Handles}
\label{subsec:morse-theory-low-dimensional}
%%%%%%%%%%%%%%%%%%%%%%%%%%%%%%%%%%%%%%%%%%%%%%%%%%%%%%%%%%%%

Morse theory connects smooth functions with handle decompositions.

Given a Morse function

\[
f:M\to\R,
\]

crossing a critical value of index \(k\) corresponds to attaching a
\(k\)-handle.

Thus

\[
\boxed{
\text{critical points}
\longleftrightarrow
\text{handles}.
}
\]

This converts differential information into topological structure.

%%%%%%%%%%%%%%%%%%%%%%%%%%%%%%%%%%%%%%%%%%%%%%%%%%%%%%%%%%%%
\subsection{Heegaard Splittings and Morse Theory}
\label{subsec:heegaard-morse}
%%%%%%%%%%%%%%%%%%%%%%%%%%%%%%%%%%%%%%%%%%%%%%%%%%%%%%%%%%%%

For a closed \(3\)-manifold, a suitable Morse function produces a Heegaard
splitting.

The lower-index handles form one handlebody, while the higher-index handles
form the other.

Therefore

\[
\boxed{
\text{Morse function}
\longrightarrow
\text{Heegaard splitting}.
}
\]

This links differential topology with \(3\)-manifold topology.

%%%%%%%%%%%%%%%%%%%%%%%%%%%%%%%%%%%%%%%%%%%%%%%%%%%%%%%%%%%%
\subsection{Knots as Boundaries}
\label{subsec:knots-as-boundaries}
%%%%%%%%%%%%%%%%%%%%%%%%%%%%%%%%%%%%%%%%%%%%%%%%%%%%%%%%%%%%

A knot can appear as the boundary of several different kinds of surfaces.

In \(S^3\),

\[
\boxed{
K=\partial F
}
\]

for a Seifert surface \(F\).

In \(D^4\),

\[
\boxed{
K=\partial F_4
}
\]

for a properly embedded surface \(F_4\).

Comparing these surfaces produces invariants of both knots and
\(4\)-manifolds.

%%%%%%%%%%%%%%%%%%%%%%%%%%%%%%%%%%%%%%%%%%%%%%%%%%%%%%%%%%%%
\subsection{Low-Dimensional Topological Invariants}
\label{subsec:low-dimensional-invariants}
%%%%%%%%%%%%%%%%%%%%%%%%%%%%%%%%%%%%%%%%%%%%%%%%%%%%%%%%%%%%

Different dimensions emphasize different invariants.

For surfaces:

\[
\boxed{
\chi,
\qquad
\text{orientability},
\qquad
g.
}
\]

For \(3\)-manifolds:

\[
\boxed{
\pi_1,
\qquad
\text{homology},
\qquad
\text{geometric decomposition},
\qquad
\text{hyperbolic volume}.
}
\]

For \(4\)-manifolds:

\[
\boxed{
Q_M,
\qquad
\sigma(M),
\qquad
\text{smooth structures},
\qquad
\text{gauge-theoretic invariants}.
}
\]

%%%%%%%%%%%%%%%%%%%%%%%%%%%%%%%%%%%%%%%%%%%%%%%%%%%%%%%%%%%%
\subsection{Floer Theory}
\label{subsec:floer-theory-low-dimensional}
%%%%%%%%%%%%%%%%%%%%%%%%%%%%%%%%%%%%%%%%%%%%%%%%%%%%%%%%%%%%

Floer theories assign algebraic structures to low-dimensional topological
objects.

Examples include

\[
\boxed{
\text{instanton Floer homology},
}
\]

\[
\boxed{
\text{monopole Floer homology},
}
\]

and

\[
\boxed{
\text{Heegaard Floer homology}.
}
\]

These theories connect

\[
3\text{-manifolds},
\qquad
4\text{-manifolds},
\qquad
\text{knots}.
\]

Their detailed construction requires substantial additional machinery and
is deferred to more advanced treatments.

%%%%%%%%%%%%%%%%%%%%%%%%%%%%%%%%%%%%%%%%%%%%%%%%%%%%%%%%%%%%
\subsection{Khovanov Homology}
\label{subsec:khovanov-low-dimensional}
%%%%%%%%%%%%%%%%%%%%%%%%%%%%%%%%%%%%%%%%%%%%%%%%%%%%%%%%%%%%

Khovanov homology assigns a bigraded homology theory to a link.

Its graded Euler characteristic recovers the Jones polynomial under standard
normalization conventions.

Thus it replaces a polynomial invariant with a richer algebraic object.

This process is an example of \emph{categorification}.

Schematically,

\[
\boxed{
\text{Jones polynomial}
\longrightarrow
\text{Khovanov homology}.
}
\]

%%%%%%%%%%%%%%%%%%%%%%%%%%%%%%%%%%%%%%%%%%%%%%%%%%%%%%%%%%%%
\subsection{Topological Quantum Field Theory}
\label{subsec:tqft-low-dimensional}
%%%%%%%%%%%%%%%%%%%%%%%%%%%%%%%%%%%%%%%%%%%%%%%%%%%%%%%%%%%%

Low-dimensional topology is deeply connected with topological quantum field
theory.

A topological quantum field theory, or TQFT, associates algebraic data to
manifolds and cobordisms.

Schematically,

\[
\boxed{
\text{manifold}
\longrightarrow
\text{algebraic object},
}
\]

while

\[
\boxed{
\text{cobordism}
\longrightarrow
\text{linear transformation}.
}
\]

These ideas connect topology, category theory, representation theory, and
mathematical physics.

%%%%%%%%%%%%%%%%%%%%%%%%%%%%%%%%%%%%%%%%%%%%%%%%%%%%%%%%%%%%
\subsection{Triangulations and Computation}
\label{subsec:low-dimensional-computation}
%%%%%%%%%%%%%%%%%%%%%%%%%%%%%%%%%%%%%%%%%%%%%%%%%%%%%%%%%%%%

Low-dimensional manifolds can often be represented combinatorially.

For surfaces, use

\[
\boxed{
\text{triangles}.
}
\]

For \(3\)-manifolds, use

\[
\boxed{
\text{tetrahedra}.
}
\]

For \(4\)-manifolds, one may use

\[
\boxed{
\text{simplices or handle decompositions}.
}
\]

These representations permit algorithmic investigation of topological
questions.

%%%%%%%%%%%%%%%%%%%%%%%%%%%%%%%%%%%%%%%%%%%%%%%%%%%%%%%%%%%%
\subsection{Normal Surface Theory}
\label{subsec:normal-surface-theory}
%%%%%%%%%%%%%%%%%%%%%%%%%%%%%%%%%%%%%%%%%%%%%%%%%%%%%%%%%%%%

Normal surface theory studies surfaces inside triangulated \(3\)-manifolds
by placing them into a standardized position relative to each tetrahedron.

Within a tetrahedron, a normal surface intersects in certain standard disk
types.

The global surface can then be encoded using integer coordinates.

Thus

\[
\boxed{
\text{embedded surface}
\longrightarrow
\text{integer vector}.
}
\]

This makes some topological questions algorithmically approachable.

%%%%%%%%%%%%%%%%%%%%%%%%%%%%%%%%%%%%%%%%%%%%%%%%%%%%%%%%%%%%
\subsection{Algorithmic Three-Manifold Topology}
\label{subsec:algorithmic-three-manifold-topology}
%%%%%%%%%%%%%%%%%%%%%%%%%%%%%%%%%%%%%%%%%%%%%%%%%%%%%%%%%%%%

Algorithms exist for important problems such as

\[
\boxed{
\text{unknot recognition},
}
\]

\[
\boxed{
3\text{-sphere recognition},
}
\]

and various manifold decomposition problems.

These algorithms combine

\[
\boxed{
\text{topology},
\quad
\text{combinatorics},
\quad
\text{algebra},
\quad
\text{computation}.
}
\]

%%%%%%%%%%%%%%%%%%%%%%%%%%%%%%%%%%%%%%%%%%%%%%%%%%%%%%%%%%%%
\subsection{A Dimensional Comparison}
\label{subsec:dimensional-comparison-topology}
%%%%%%%%%%%%%%%%%%%%%%%%%%%%%%%%%%%%%%%%%%%%%%%%%%%%%%%%%%%%

A useful comparison is

\[
\boxed{
\begin{array}{c|c}
\text{Dimension} & \text{Characteristic Themes}\\
\hline
2 & \text{classification by genus and orientability}\\
3 & \text{knots, decompositions, geometry, hyperbolicity}\\
4 & \text{intersection forms, gauge theory, exotic smoothness}
\end{array}
}
\]

This does not mean higher dimensions are simple.

Rather, dimensions \(2\), \(3\), and \(4\) possess special techniques and
phenomena that justify treating them separately.

%%%%%%%%%%%%%%%%%%%%%%%%%%%%%%%%%%%%%%%%%%%%%%%%%%%%%%%%%%%%
\subsection{Worked Example: Classifying a Surface}
\label{subsec:worked-classifying-surface}
%%%%%%%%%%%%%%%%%%%%%%%%%%%%%%%%%%%%%%%%%%%%%%%%%%%%%%%%%%%%

\begin{workedexamplebox}{Orientable Surface Classification}

Suppose \(M\) is a closed connected orientable surface with

\[
\chi(M)=-4.
\]

For an orientable surface,

\[
\chi=2-2g.
\]

Therefore

\[
-4=2-2g.
\]

Hence

\[
-6=-2g,
\]

so

\[
g=3.
\]

By the classification theorem,

\[
M\cong\Sigma_3.
\]

\begin{answerbox}
\[
\boxed{
M\cong T^2\#T^2\#T^2.
}
\]
\end{answerbox}
\end{workedexamplebox}

%%%%%%%%%%%%%%%%%%%%%%%%%%%%%%%%%%%%%%%%%%%%%%%%%%%%%%%%%%%%
\subsection{Worked Example: Nonorientable Classification}
\label{subsec:worked-nonorientable-classification}
%%%%%%%%%%%%%%%%%%%%%%%%%%%%%%%%%%%%%%%%%%%%%%%%%%%%%%%%%%%%

\begin{workedexamplebox}{Finding the Crosscap Number}

Suppose \(N\) is a closed connected nonorientable surface with

\[
\chi(N)=-3.
\]

For a nonorientable surface,

\[
\chi=2-k.
\]

Therefore

\[
-3=2-k.
\]

Hence

\[
k=5.
\]

Thus

\[
N
\cong
\#^5\mathbb{RP}^2.
\]

\begin{answerbox}
\[
\boxed{
N\cong
\mathbb{RP}^2
\#
\mathbb{RP}^2
\#
\mathbb{RP}^2
\#
\mathbb{RP}^2
\#
\mathbb{RP}^2.
}
\]
\end{answerbox}
\end{workedexamplebox}

%%%%%%%%%%%%%%%%%%%%%%%%%%%%%%%%%%%%%%%%%%%%%%%%%%%%%%%%%%%%
\subsection{Worked Example: Signature}
\label{subsec:worked-four-manifold-signature}
%%%%%%%%%%%%%%%%%%%%%%%%%%%%%%%%%%%%%%%%%%%%%%%%%%%%%%%%%%%%

\begin{workedexamplebox}{Computing the Signature}

Suppose the intersection form of a \(4\)-manifold is represented over
\(\R\) by

\[
Q=
\begin{pmatrix}
1&0&0&0\\
0&1&0&0\\
0&0&-1&0\\
0&0&0&-1
\end{pmatrix}.
\]

There are two positive eigenvalues and two negative eigenvalues.

Therefore

\[
b_2^+=2
\]

and

\[
b_2^-=2.
\]

Hence

\[
\sigma(M)
=
b_2^+-b_2^-
=
2-2.
\]

\begin{answerbox}
\[
\boxed{
\sigma(M)=0.
}
\]
\end{answerbox}
\end{workedexamplebox}

%%%%%%%%%%%%%%%%%%%%%%%%%%%%%%%%%%%%%%%%%%%%%%%%%%%%%%%%%%%%
\subsection{Common Errors}
\label{subsec:low-dimensional-common-errors}
%%%%%%%%%%%%%%%%%%%%%%%%%%%%%%%%%%%%%%%%%%%%%%%%%%%%%%%%%%%%

\subsubsection{Assuming Euler Characteristic Alone Classifies All Surfaces}

For closed connected surfaces, one must also know orientability.

For example,

\[
T^2
\]

and the Klein bottle both satisfy

\[
\chi=0,
\]

but they are not homeomorphic.

%%%%%%%%%%%%%%%%%%%%%%%%%%%%%%%%%%%%%%%%%%%%%%%%%%%%%%%%%%%%
\subsubsection{Confusing Genus and Euler Characteristic}

For a closed orientable surface,

\[
\chi=2-2g.
\]

The quantities are related but are not identical.

%%%%%%%%%%%%%%%%%%%%%%%%%%%%%%%%%%%%%%%%%%%%%%%%%%%%%%%%%%%%
\subsubsection{Using the Orientable Formula for a Nonorientable Surface}

For orientable surfaces,

\[
\chi=2-2g-b.
\]

For nonorientable surfaces,

\[
\chi=2-k-b.
\]

%%%%%%%%%%%%%%%%%%%%%%%%%%%%%%%%%%%%%%%%%%%%%%%%%%%%%%%%%%%%
\subsubsection{Assuming the Torus and Klein Bottle Are Equivalent}

Both have Euler characteristic zero, but the torus is orientable while the
Klein bottle is not.

%%%%%%%%%%%%%%%%%%%%%%%%%%%%%%%%%%%%%%%%%%%%%%%%%%%%%%%%%%%%
\subsubsection{Confusing a Handlebody with Its Boundary}

A genus-\(g\) handlebody is three-dimensional.

Its boundary

\[
\Sigma_g
\]

is two-dimensional.

%%%%%%%%%%%%%%%%%%%%%%%%%%%%%%%%%%%%%%%%%%%%%%%%%%%%%%%%%%%%
\subsubsection{Assuming Every Prime Three-Manifold Is Irreducible}

The standard orientable exception is

\[
\boxed{
S^2\times S^1.
}
\]

%%%%%%%%%%%%%%%%%%%%%%%%%%%%%%%%%%%%%%%%%%%%%%%%%%%%%%%%%%%%
\subsubsection{Assuming Every Three-Manifold Is Hyperbolic}

The eight Thurston geometries include several nonhyperbolic possibilities.

Additionally, a manifold may need to be decomposed into geometric pieces.

%%%%%%%%%%%%%%%%%%%%%%%%%%%%%%%%%%%%%%%%%%%%%%%%%%%%%%%%%%%%
\subsubsection{Confusing Poincar\'e Conjecture with Higher-Dimensional Sphere Recognition}

The classical Poincar\'e conjecture concerns closed simply connected
\(3\)-manifolds.

Analogous questions in other dimensions have distinct histories and
techniques.

%%%%%%%%%%%%%%%%%%%%%%%%%%%%%%%%%%%%%%%%%%%%%%%%%%%%%%%%%%%%
\subsubsection{Assuming Homeomorphic Smooth Manifolds Are Diffeomorphic}

This fails dramatically in dimension four.

%%%%%%%%%%%%%%%%%%%%%%%%%%%%%%%%%%%%%%%%%%%%%%%%%%%%%%%%%%%%
\subsubsection{Confusing Knot Genus and Slice Genus}

The knot genus uses surfaces in

\[
S^3,
\]

whereas the slice genus uses surfaces in

\[
D^4.
\]

In general,

\[
g_4(K)\leq g(K).
\]

%%%%%%%%%%%%%%%%%%%%%%%%%%%%%%%%%%%%%%%%%%%%%%%%%%%%%%%%%%%%
\subsubsection{Assuming Every Topological Four-Manifold Has a Unique Smooth Structure}

Dimension four admits exotic smooth structures.

%%%%%%%%%%%%%%%%%%%%%%%%%%%%%%%%%%%%%%%%%%%%%%%%%%%%%%%%%%%%
\subsubsection{Thinking Kirby Diagrams Are Ordinary Knot Diagrams}

Kirby diagrams contain additional handle and framing information.

They encode \(4\)-manifolds rather than merely knots.

%%%%%%%%%%%%%%%%%%%%%%%%%%%%%%%%%%%%%%%%%%%%%%%%%%%%%%%%%%%%
\subsection{Applications and Connections}
\label{subsec:low-dimensional-applications}
%%%%%%%%%%%%%%%%%%%%%%%%%%%%%%%%%%%%%%%%%%%%%%%%%%%%%%%%%%%%

\begin{applicationbox}
\textbf{Surface Theory.}
The classification theorem provides a complete description of compact
connected surfaces through orientability, genus, and boundary components.

\medskip

\textbf{Algebraic Topology.}
Fundamental groups, homology, cohomology, covering spaces, and intersection
forms provide algebraic invariants of low-dimensional manifolds.

\medskip

\textbf{Knot Theory.}
Knots live in \(3\)-manifolds, their complements are \(3\)-manifolds, and
slice knots connect classical knot theory with \(4\)-manifolds.

\medskip

\textbf{Hyperbolic Geometry.}
Many \(3\)-manifolds and knot complements carry hyperbolic structures.

\medskip

\textbf{Differential Geometry.}
Ricci flow connects curvature with the topology and geometry of
\(3\)-manifolds.

\medskip

\textbf{Group Theory.}
Surface groups, mapping class groups, knot groups, and \(3\)-manifold groups
provide central examples of infinite groups.

\medskip

\textbf{Linear Algebra.}
Intersection forms convert geometric intersections of surfaces in
\(4\)-manifolds into symmetric bilinear forms.

\medskip

\textbf{Partial Differential Equations.}
Ricci flow is a nonlinear geometric PDE whose analysis led to the proof of
geometrization.

\medskip

\textbf{Mathematical Physics.}
Gauge theory, Floer theory, quantum topology, and TQFT create major
connections between low-dimensional topology and physics.

\medskip

\textbf{Computation.}
Triangulations and normal surface theory turn some manifold problems into
finite combinatorial and algorithmic problems.
\end{applicationbox}

%%%%%%%%%%%%%%%%%%%%%%%%%%%%%%%%%%%%%%%%%%%%%%%%%%%%%%%%%%%%
\subsection{Exercises}
\label{subsec:low-dimensional-exercises}
%%%%%%%%%%%%%%%%%%%%%%%%%%%%%%%%%%%%%%%%%%%%%%%%%%%%%%%%%%%%

\begin{exercise}[1]
What dimensions are usually emphasized in low-dimensional topology?
\end{exercise}

\begin{exercise}[1]
Define a surface.
\end{exercise}

\begin{exercise}[1]
State the Euler characteristic of \(S^2\).
\end{exercise}

\begin{exercise}[1]
State the Euler characteristic of \(T^2\).
\end{exercise}

\begin{exercise}[1]
Is \(\mathbb{RP}^2\) orientable?
\end{exercise}

\begin{exercise}[1]
Is the Klein bottle orientable?
\end{exercise}

\begin{exercise}[1]
Define connected sum.
\end{exercise}

\begin{exercise}[1]
What is the genus of the torus?
\end{exercise}

\begin{exercise}[1]
What is a Heegaard splitting?
\end{exercise}

\begin{exercise}[1]
What is an exotic smooth structure?
\end{exercise}

\begin{exercise}[2]
Compute the Euler characteristic of a closed orientable surface of genus
\(5\).
\end{exercise}

\begin{exercise}[2]
A closed orientable surface has Euler characteristic \(-8\). Determine its
genus.
\end{exercise}

\begin{exercise}[2]
Compute the Euler characteristic of a nonorientable surface with \(7\)
crosscaps.
\end{exercise}

\begin{exercise}[2]
A closed nonorientable surface has Euler characteristic \(-4\). Determine
its crosscap number.
\end{exercise}

\begin{exercise}[2]
Compute the Euler characteristic of an orientable genus-\(3\) surface with
\(4\) boundary components.
\end{exercise}

\begin{exercise}[2]
Compute the Euler characteristic of a nonorientable surface with \(5\)
crosscaps and \(2\) boundary components.
\end{exercise}

\begin{exercise}[2]
Explain why the torus and Klein bottle are not homeomorphic even though both
have Euler characteristic zero.
\end{exercise}

\begin{exercise}[2]
Compute

\[
\pi_1(S^2).
\]
\end{exercise}

\begin{exercise}[2]
Compute

\[
\pi_1(T^2).
\]
\end{exercise}

\begin{exercise}[2]
Compute

\[
\pi_1(\mathbb{RP}^2).
\]
\end{exercise}

\begin{exercise}[2]
Describe the polygonal word for the torus.
\end{exercise}

\begin{exercise}[2]
Describe the polygonal word for the Klein bottle.
\end{exercise}

\begin{exercise}[2]
What is the Heegaard genus of \(S^3\)?
\end{exercise}

\begin{exercise}[2]
What is the boundary of a genus-\(g\) handlebody?
\end{exercise}

\begin{exercise}[2]
List the eight Thurston geometries.
\end{exercise}

\begin{exercise}[2]
State the Poincar\'e conjecture.
\end{exercise}

\begin{exercise}[2]
Explain the difference between homeomorphic and diffeomorphic manifolds.
\end{exercise}

\begin{exercise}[2]
Compute the signature of

\[
Q=
\begin{pmatrix}
1&0&0\\
0&-1&0\\
0&0&-1
\end{pmatrix}.
\]
\end{exercise}

\begin{exercise}[2]
If a knot has

\[
g(K)=4,
\]

what can be concluded about the possible values of \(g_4(K)\)?
\end{exercise}

\begin{exercise}[3]
Prove that

\[
\chi(\Sigma_g)=2-2g.
\]
\end{exercise}

\begin{exercise}[3]
Show that

\[
\chi(M\#N)
=
\chi(M)+\chi(N)-2
\]

for closed connected surfaces.
\end{exercise}

\begin{exercise}[3]
Derive a presentation of

\[
\pi_1(\Sigma_g).
\]
\end{exercise}

\begin{exercise}[3]
Show that the abelianization of

\[
\pi_1(\Sigma_g)
\]

is

\[
\Z^{2g}.
\]
\end{exercise}

\begin{exercise}[3]
Classify every closed connected surface with Euler characteristic zero.
\end{exercise}

\begin{exercise}[3]
Explain why orientability and Euler characteristic classify closed connected
surfaces.
\end{exercise}

\begin{exercise}[3]
Describe a Dehn twist geometrically.
\end{exercise}

\begin{exercise}[3]
Explain why mapping class groups arise naturally in Heegaard splittings.
\end{exercise}

\begin{exercise}[3]
Describe a genus-\(1\) Heegaard splitting.
\end{exercise}

\begin{exercise}[3]
Explain why knot complements are natural objects in \(3\)-manifold topology.
\end{exercise}

\begin{exercise}[3]
Explain the difference between prime and irreducible \(3\)-manifolds.
\end{exercise}

\begin{exercise}[3]
Describe the purpose of the JSJ decomposition.
\end{exercise}

\begin{exercise}[3]
Explain why hyperbolic volume becomes a topological invariant for
finite-volume hyperbolic \(3\)-manifolds.
\end{exercise}

\begin{exercise}[3]
Explain how Dehn surgery on a framed link produces a \(3\)-manifold.
\end{exercise}

\begin{exercise}[3]
For

\[
Q=
\begin{pmatrix}
1&0&0&0&0\\
0&1&0&0&0\\
0&0&1&0&0\\
0&0&0&-1&0\\
0&0&0&0&-1
\end{pmatrix},
\]

compute \(b_2^+\), \(b_2^-\), and the signature.
\end{exercise}

\begin{exercise}[3]
Explain why

\[
g_4(K)\leq g(K).
\]
\end{exercise}

\begin{exercise}[3]
Describe how a framed knot can encode a \(2\)-handle attachment.
\end{exercise}

\begin{exercise}[3]
Explain the relationship between Morse critical points and handles.
\end{exercise}

\begin{exercise}[4]
Study a proof of the classification theorem for compact connected surfaces.
\end{exercise}

\begin{exercise}[4]
Research the mapping class group of the torus and its relationship with

\[
\mathrm{SL}_2(\Z).
\]
\end{exercise}

\begin{exercise}[4]
Study the theorem that Dehn twists generate mapping class groups of compact
orientable surfaces.
\end{exercise}

\begin{exercise}[4]
Research the prime decomposition theorem for closed orientable
\(3\)-manifolds.
\end{exercise}

\begin{exercise}[4]
Study normal surface theory and explain how embedded surfaces can be encoded
by integer vectors.
\end{exercise}

\begin{exercise}[4]
Research the JSJ decomposition of \(3\)-manifolds.
\end{exercise}

\begin{exercise}[4]
Study the figure-eight knot complement as a hyperbolic \(3\)-manifold.
\end{exercise}

\begin{exercise}[4]
Research the role of Ricci flow in the proof of geometrization.
\end{exercise}

\begin{exercise}[4]
Study the Lickorish--Wallace theorem.
\end{exercise}

\begin{exercise}[4]
Research the classification of lens spaces.
\end{exercise}

\begin{exercise}[4]
Study Freedman's classification results for simply connected topological
\(4\)-manifolds.
\end{exercise}

\begin{exercise}[4]
Study Donaldson's diagonalization theorem.
\end{exercise}

\begin{exercise}[4]
Explain how the combination of Freedman's and Donaldson's results reveals
the exceptional behavior of dimension four.
\end{exercise}

\begin{exercise}[4]
Research exotic \(\R^4\)'s.
\end{exercise}

\begin{exercise}[4]
Study Kirby calculus and compare Kirby moves with Reidemeister moves.
\end{exercise}

\begin{exercise}[4]
Research the smooth knot concordance group.
\end{exercise}

\begin{exercise}[4]
Study Heegaard Floer homology and describe one way it produces invariants of
knots or \(3\)-manifolds.
\end{exercise}

\begin{exercise}[4]
Research Khovanov homology and explain the meaning of categorification.
\end{exercise}

%%%%%%%%%%%%%%%%%%%%%%%%%%%%%%%%%%%%%%%%%%%%%%%%%%%%%%%%%%%%
\subsection{Conceptual Summary}
\label{subsec:low-dimensional-summary}
%%%%%%%%%%%%%%%%%%%%%%%%%%%%%%%%%%%%%%%%%%%%%%%%%%%%%%%%%%%%

Low-dimensional topology focuses on

\[
\boxed{
2\text{-manifolds},
\qquad
3\text{-manifolds},
\qquad
4\text{-manifolds}.
}
\]

Closed connected surfaces are completely classified.

Orientable surfaces have the form

\[
\boxed{
\Sigma_g=\#^gT^2
}
\]

with

\[
\boxed{
\chi(\Sigma_g)=2-2g.
}
\]

Nonorientable surfaces have the form

\[
\boxed{
N_k=\#^k\mathbb{RP}^2
}
\]

with

\[
\boxed{
\chi(N_k)=2-k.
}
\]

Surface topology leads naturally to mapping class groups and Dehn twists.

Three-manifolds can be studied using

\[
\boxed{
\text{Heegaard splittings},
\quad
\text{triangulations},
\quad
\text{essential surfaces},
\quad
\text{surgery}.
}
\]

Prime and JSJ decompositions break complicated \(3\)-manifolds into simpler
pieces.

Thurston's geometrization program relates these pieces to eight model
geometries.

The Poincar\'e conjecture states

\[
\boxed{
\pi_1(M)=\{e\},
\quad
M\text{ closed }3\text{-manifold}
\Longrightarrow
M\cong S^3.
}
\]

Dimension four introduces intersection forms

\[
\boxed{
Q_M:
H_2(M;\Z)/\operatorname{Tor}
\times
H_2(M;\Z)/\operatorname{Tor}
\to\Z,
}
\]

as well as exotic smooth structures.

Freedman's and Donaldson's theories reveal a major distinction between
topological and smooth \(4\)-manifolds.

Knots connect dimensions three and four through Seifert surfaces,
concordance, and slice genus:

\[
\boxed{
g_4(K)\leq g(K).
}
\]

Handle decompositions and Kirby calculus convert smooth \(4\)-manifold
questions into diagrammatic problems.

Low-dimensional topology therefore sits at the intersection of

\[
\boxed{
\text{topology},
\quad
\text{geometry},
\quad
\text{algebra},
\quad
\text{analysis},
\quad
\text{mathematical physics}.
}
\]

%%%%%%%%%%%%%%%%%%%%%%%%%%%%%%%%%%%%%%%%%%%%%%%%%%%%%%%%%%%%
\subsection{Section Connections}
\label{subsec:low-dimensional-connections}
%%%%%%%%%%%%%%%%%%%%%%%%%%%%%%%%%%%%%%%%%%%%%%%%%%%%%%%%%%%%

\chapterconnection{
Low-Dimensional Topology completes the topology chapter by combining the
ideas developed throughout Point-Set Topology, Algebraic Topology,
Differential Topology, Geometric Topology, Manifold Theory, and Knot Theory.
Surface classification provides a complete model of a manifold
classification problem. Mapping class groups connect surfaces with Group
Theory. Heegaard splittings and Dehn surgery connect surfaces and knots with
\(3\)-manifolds. Geometrization links topology with Euclidean, spherical,
and hyperbolic geometry. Ricci flow connects topology with Differential
Geometry and Partial Differential Equations. Four-manifold topology uses
homology, intersection forms, gauge theory, handle decompositions, and knot
concordance. These ideas provide a natural transition from Topology into
Chapter 5, Analysis, where continuity, convergence, metric structure,
differentiation, integration, and infinite processes are developed in
greater analytic depth.
}

%%%%%%%%%%%%%%%%%%%%%%%%%%%%%%%%%%%%%%%%%%%%%%%%%%%%%%%%%%%%
% SECTION SOURCE TRACKING
%%%%%%%%%%%%%%%%%%%%%%%%%%%%%%%%%%%%%%%%%%%%%%%%%%%%%%%%%%%%

%------------------------------------------------------------
% 4.7 Low-Dimensional Topology
%------------------------------------------------------------
%
% Source basis:
%   - Standard surface topology, 3-manifold topology,
%     4-manifold topology, geometric topology, and
%     introductory low-dimensional topology.
%
% Used for:
%   - special behavior of dimensions 2, 3, and 4;
%   - surfaces;
%   - sphere, torus, projective plane, Klein bottle;
%   - polygonal surface models;
%   - orientability;
%   - connected sums;
%   - orientable and nonorientable genus;
%   - Euler characteristic formulas;
%   - classification of compact connected surfaces;
%   - surfaces with boundary;
%   - fundamental groups of surfaces;
%   - curves and geometric intersection;
%   - mapping class groups;
%   - Dehn twists;
%   - 3-manifolds;
%   - lens spaces;
%   - handlebodies;
%   - Heegaard splittings and Heegaard genus;
%   - triangulations;
%   - prime decomposition;
%   - irreducibility;
%   - incompressible surfaces;
%   - Haken manifolds;
%   - JSJ decomposition;
%   - Thurston geometries;
%   - geometrization;
%   - Poincare conjecture;
%   - Ricci flow;
%   - Dehn surgery;
%   - surgery presentations;
%   - Lickorish-Wallace theorem;
%   - 4-manifolds;
%   - topological versus smooth structures;
%   - exotic smooth structures;
%   - exotic R^4;
%   - intersection forms;
%   - signature;
%   - definite and indefinite forms;
%   - Freedman theory;
%   - Donaldson theory;
%   - handles and Kirby diagrams;
%   - Kirby calculus;
%   - knots in 4-manifolds;
%   - slice genus;
%   - concordance;
%   - cobordism;
%   - Morse theory;
%   - Floer theories;
%   - Khovanov homology;
%   - TQFT;
%   - normal surface theory;
%   - algorithmic topology.
%
% Notes:
%   - Point-set foundations are developed in Section 4.1.
%   - Fundamental groups, homology, covering spaces, and
%     related invariants are developed in Algebraic Topology.
%   - Smooth manifolds, transversality, Morse theory, and
%     related foundations are developed in Differential
%     Topology and Manifold Theory.
%   - Ambient isotopy, manifold decomposition, and surgery
%     are connected with Geometric Topology.
%   - Knot complements, Dehn surgery, concordance, and slice
%     knots build directly on Knot Theory.
%   - Hyperbolic geometry is developed in the Geometry
%     chapter.
%   - Ricci flow and curvature connect with Differential
%     Geometry and later Differential Equations.
%   - Freedman theory, Donaldson theory, Floer homology,
%     Khovanov homology, gauge theory, and TQFT are introduced
%     structurally rather than developed in full.
%
%%%%%%%%%%%%%%%%%%%%%%%%%%%%%%%%%%%%%%%%%%%%%%%%%%%%%%%%%%%%